% English translation of questions/5780/semA-moedA/partA-q2.tex (metadata is taken from the Hebrew file).

\begin{question}
Find the largest value and the smallest value of the function $y=\sqrt[3]{x}(1+2x)$ on the interval $[-8,1]$.
\end{question}

\begin{hint}
The candidates are the endpoints of the interval, points where $y'=0$, and points where $y$ is not differentiable (pay attention to $x=0$).
\end{hint}

\begin{solution}
The function $y=x^{1/3}+2x^{4/3}$ is continuous on the closed interval $[-8,1]$, hence by \textbf{Weierstrass's theorem} it attains a maximum and a minimum there. These points are at the endpoints of the interval or at interior critical points (where $y'=0$ or $y'$ does not exist), by Fermat's theorem.

\textbf{Derivative} (for $x\ne0$):
\[ y' = \frac13x^{-2/3}+\frac83x^{1/3} = \frac{1+8x}{3\sqrt[3]{x^2}}. \]
\begin{itemize}
  \item $y'=0\iff x=-\frac18\in(-8,1)$.
  \item At $x=0$ the derivative does not exist: $\frac{y(h)-y(0)}{h}=\frac{h^{1/3}(1+2h)}{h}=\frac{1+2h}{h^{2/3}}\to+\infty$. Hence $x=0$ is a critical point.
\end{itemize}

\textbf{Values of the function at the candidates:}
\begin{align*}
y(-8) &= \sqrt[3]{-8}\,(1-16) = (-2)(-15) = 30,\\
y\left(-\tfrac18\right) &= \sqrt[3]{-\tfrac18}\left(1-\tfrac14\right) = -\tfrac12\cdot\tfrac34 = -\tfrac38,\\
y(0) &= 0,\\
y(1) &= 1\cdot 3 = 3.
\end{align*}

\textbf{Answer:} the largest value is $30$ (attained at $x=-8$), and the smallest value is $-\frac38$ (attained at $x=-\frac18$).
\end{solution}
